\section{Introduction}

This paper expands on LP-208 (\texttt{lps/LP-208-dag-mempool.md})
with detail on the wire format for the two new ZAP schemas
(\texttt{0xF0}~DAGHeader, \texttt{0xF1}~DAGBody), the parent-linking
discipline, the throughput model, the garbage collection contract,
and the composition with sibling LPs.

\subsection{What was wrong before}

Today's Lux mempool model: each validator maintains its own ranked
heap of pending transactions; a leader at each round picks from its
local mempool when proposing a block. Aggregate throughput is
bounded by the slowest validator -- if one validator's mempool can
admit \SI{50000}{tx/s}, that is the aggregate ceiling no matter how
many validators participate.

Three concrete failures:

\begin{enumerate}[leftmargin=2em,nosep]
\item \textbf{Throughput is min-bounded.} Aggregate admission rate at
  a chain is $\min(\textrm{per-validator-admission})$. For a 64-validator
  chain with one validator on consumer hardware
  (\SI{10000}{tx/s}) and 63 on Blackwell (\SI{200000}{tx/s}),
  aggregate is \SI{10000}{tx/s} -- $99.2\%$ of the validator-set
  capacity is wasted.
\item \textbf{Tx loss on leader byzantine.} A leader at round $R$
  picks from its local mempool. If the leader is byzantine and
  withholds its block, all txs in its local mempool are stuck until
  the next round; honest validators' versions of those txs may have
  already expired.
\item \textbf{Latency is leader-rotation-dependent.} A tx admitted
  to validator $V$'s mempool at time $T$ does not propagate into a
  block until $V$ is the leader at some future round. The gap is
  bounded by leader-rotation period; for round-robin with $N$
  validators and block period $P$, the worst-case tx-to-block
  latency is $N \times P$.
\end{enumerate}

\subsection{What this LP solves}

The DAG mempool inverts the model. Every validator continuously
broadcasts a stream of ZAP headers; each header carries a list of
tx hashes, parent pointers to the most recent headers from peer
validators, and a leader-signed digest. The actual tx bytes live in
the Kademlia DHT (LP-201 layer~C), content-addressed by hash. The
aggregate structure is a directed acyclic graph: every header has
parents from peers, so the graph captures the partial order of tx
admission across the validator set.

Throughput becomes $\sum(\textrm{validators})$. If each validator
emits \SI{50000}{tx-hashes/s}, a 64-validator set admits
\SI{3.2}{\mega tx-hashes/s} into the DAG. The bottleneck moves from
mempool admission to DHT throughput and signature-verification rate
per validator -- both linearly scalable.

\subsection{Scope of this LP}

DAG total-ordering -- the step that picks a single canonical
sequence from the DAG for execution -- is deferred to LP-209
(Mysticeti, \cite{mysticeti2024}). This LP defines only the
construction surface: the wire format for headers and bodies, the
parent-linking discipline, the DHT-storage contract for bodies, the
garbage-collection rule for old headers, and the per-validator
emission rate.

This separation is deliberate. Multiple ordering algorithms can
consume the same DAG: Mysticeti (committee-based ordering with
sub-second finality); Bullshark \cite{bullshark2022} (round-robin
anchored ordering with FIFO fairness); Tusk (asynchronous ordering
with worst-case finality). Each algorithm has a different
latency/throughput/fairness tradeoff. The DAG construction is
invariant across these choices.
